\begin{questionS}
A particular synchronisation object is to be used by two threads, which we
call the \emph{parent} and the \emph{child}.  The child repeatedly calls an
operation |signalUpAndWait|, while the parent alternates between calling
operations {\scalastyle wait}\-|For|\-|Signal|\-|Up| and |signalDown|.  The
|waitForSignalUp| operation blocks until the child calls
|signal|\-|Up|\-|And|\-|Wait|.  The |signalUpAndWait| operation blocks until
the parent calls |signal|\-|Down|.  Thus such an object can be used to
repeatedly send two signals, one in each direction.  Implement a class of such
objects with the following signature.
%
\begin{scala}
class Signal{
  /** Signal to the parent, and wait for a signal back. */
  def signalUpAndWait(): Unit = ...

  /** Wait for a signal from the child. */
  def waitForSignalUp(): Unit = ...

  /** Signal to the child. */
  def signalDown(): Unit = ...
} 
\end{scala}

Now implement a barrier synchronisation object that uses multiple |Signal|
objects internally.  Each synchronisation of $n$~threads should run in
time~$O(n)$.
\end{questionS}

% =======================================================

\begin{answerS}
In the |Signal| objects, it is enough to maintain a boolean that is true if the
child has signalled to the parent, but not yet received a signal back.  Each
signal can then be straightforwardly implemented using a |wait| and a |notify|.
%
\begin{scala}
class Signal{
  /** The state of this object.  £true£ represents that the child has signalled,
    * but not yet received a signal back. */
  private var state = false

  /** Signal to the parent, and wait for a signal back. */
  def signalUpAndWait() = synchronized{
    require(!state); state = true; notify()
    while(state) wait()
  }

  /** Wait for a signal from the child. */
  def waitForSignalUp() = synchronized{ while(!state) wait() }

  /** Signal to the child. */
  def signalDown() = synchronized{ state = false; notify() }
}
\end{scala}

We build the barrier object by arranging threads into a heap, with a |Signal|
object between each child and parent.  The implementation is then very similar
to that for Exercise~\ref{ex:barrierLog}, except replacing channel
communications with operations on the |Signal|s.
%
\begin{scala}
class HeapBarrier(n: Int){
  /** The objects used for signalling.  £signals(i)£ is used for signalling between 
    * thread £i£ and its parent. */
  private val signals = Array.fill(n)(new Signal)

  /** Perform a barrier synchronisation.
    * @param me the unique identity of this thread. */
  def sync(me: Int) = {
    require(0 <= me && me < n)
    val child1 = 2*me+1; val child2 = 2*me+2
    // Wait for children
    if(child1 < n) signals(child1).waitForSignalUp()
    if(child2 < n) signals(child2).waitForSignalUp()
    // Signal to parent and wait for signal back
    if(me != 0) signals(me).signalUpAndWait()
    // Signal to children
    if(child1 < n) signals(child1).signalDown()
    if(child2 < n) signals(child2).signalDown()
  }
}
\end{scala}

This is basically the same as the implementation of barriers in SCL.  (The SCL
implementation gives better feedback when one of the preconditions fails; this
indicates that the client code isn't using thread identities correctly.)  The
operations on |Signal|s are lighter weight than channel communications, so
this version is likely to be more efficient. 

We could adapt this to implement a combining barrier by arranging for each
signal to pass a piece of data, by giving operations the following types.
%
\begin{scala}
class Signal[A]{
  def signalUpAndWait(x: A): A = ...
  def waitForSignalUp(): A = ...
  def signalDown(): A = ...
} 
\end{scala}
\end{answerS}
